\scanpage{23}

Prerequisites for the courses I and II

I) Introduction to A.G. : Category theory, basic algebra including some commutative algebra. Bourbaki is a good reference. Most facts we are using will be recalled.

II) Introduction to algebraic groups. Familiarity with the basic notions of A.G. will be assumed.

A.G. is a messy subject, \inserted{\unsure{even upon foundations|even more foundations},} we cannot give anything like a detailed exposition of all of it in one summer course. We will have to give rather a survey of some key points, being rather sketchy for proofs, -- most of which are straightforward. We can \unsure{pile in|pile up|pick up|fill in} whatever is requested. Homework will be necessary for those who do not have already a fair background.

Question is how to \unsure{relate|belt} the two courses. The contents of \deleted{\unsure{II}} \unsure{I are|are} prerequisites for II !

\textcircled{1}\quad $k$

\[
P_I=P[(T_i)_{i\in I}]
\]

\[
k^I\simeq\mathrm{Hom}_{k\text{-}\mathrm{alg}}(P_I,k)
\qquad\inserted{\text{l'espace affine sur $k$ de dim $I$}}\quad\deleted{\unsure{V_{\mathrm{alg}}(P_I,k)}}
\]

\margin{$I,J$ finis ou infinis}

\[
\begin{aligned}
&f\in P_I,\quad x=(a_i)_{i\in I}\in k^I\qquad\deleted{\unsure{f\in k^{\prime I}}}\\
&f(x)=\varepsilon_x(f)
\end{aligned}
\]

\[
F_j(t)=0\qquad j\in J\qquad V((F_j)) \tag*{(1)}
\]

\textcircled{2} (a) Pbs \inserted{du type} arithmétiques : très \unsure{différents|affines|différent} suivant la nature de $k$ : \ldots, $\mathbb R,\mathbb C$

\textcircled{3} (b) Pbs \inserted{de type} géométriques, issu de la géométrie sur $\mathbb R$ \deleted{et} sur $\mathbb C$ on cherche les solutions de (1) \deleted{pour} \uline{toute} \deleted{extension} \inserted{$k$-algèbre} $k'$, \deleted{\unsure{de $k$}} et comment elles varient ; et les propriétés de (1) qui sont stables par de tels changements \ldots

\[
\begin{aligned}
\mathbb E_k^I:\quad k'&\longmapsto k^{\prime I} &\mathbb E_k^I(k')&=k^{\prime I}\\
V_S(k')&\subset\mathbb E_k^I(k') & V_S&\subset\mathbb E_k^I
\left\{\begin{aligned}&\text{étude indépendante du plongement}\\ &\text{étude \unsure{du plongement|des plongements}}\end{aligned}\right.
\end{aligned}
\]

% archive sheet number: 22
